\section{Discussion}
\label{sec:discussion}

\textbf{What the four reversals have in common.} Each is a ratio whose denominator was never established: a contrast with no null, a recall with no positive-prediction rate, a score over an unknown number of label slots, a comparison across two output regimes. The fix is procedural, not architectural. All four audits run on artifacts a benchmark already writes, on CPU, in minutes.

\textbf{A reporting checklist.} We propose that any claim of the form ``model $M$ relies on the image,'' ``model $M$ is robust,'' or ``$M$ transfers'' be accompanied by: (1) output-support size and top-share on the real condition; (2) a real-vs-real null for the contrast; (3) positive recall, positive-prediction rate, and precision; (4) contract coverage and output-mode indicators; (5) the effective input resolution and precision; and (6) denominators for every metric. The first four are the preconditions; the last two are measurement declarations.

\textbf{Why the literature needs re-reading.} Prior no-image ablations report contrasts without a null and without an output-support check. Our results suggest at least two ways those numbers can mislead: a refusal template produces a large contrast with no evidence extraction, and a collapsed model produces a contrast of zero for lack of variation. Neither is a criticism of any specific paper's execution; it is a statement that the statistic as commonly computed does not identify the target quantity.

\textbf{Silent failures are the deployment-relevant quadrant.} If a monitoring system watches for uncertainty, refusal, or reasoning-channel activation, it will catch resolution loss and blur and miss noise. This asymmetry is observed in MedGemma-1.5, the only model in our matrix with a reasoning channel we can read: its marker rate is 71.98\% for resolution and 0.46\% for noise while the recall damage is comparable. The other five models do not expose this channel, so this is single-model evidence, not a cross-model law. We are also careful about what it does and does not show. The marker rate is a \emph{model-side} signal, not a measure of human or image-side visibility; a no-reference Laplacian-variance statistic separates clean from noise-severe at AUC $1.000$ on the same data, and a rule-based search finds no explicit quality language in any noise-severe reply (Appendix~\ref{app:risk}, Table~\ref{tab:imageside}, and \url{scripts/83_quality_language_audit.py}). The claim is therefore narrower: \emph{this model's own monitoring and reporting signals do not detect a corruption that demonstrably damages its answers}. A practical monitor should combine an image-side corruption detector with a model-side hesitation signal, because the two fail on opposite ends of the corruption spectrum; whether the model-side result generalizes requires a model with a comparable readable channel.

\textbf{Contract is part of the model's observable behaviour.} The schema ablation shows that the same weights can be a confident-normal shortcut or an abstaining one depending on the requested output contract. This is both a measurement warning (the contract is a variable) and a design lever (if abstention is desired, it can be requested). It also means ``the model is overconfident on blank images'' is not a model-level fact.

\textbf{What we could not establish.} The selective-trust framework that motivated this study is not validated: its components change error ranking only marginally and do not produce coverage. We report that as a negative result rather than fitting a more permissive threshold after seeing the test set.

